% !TeX root = ../../main.tex
% ============================================================
%  CHAPTER 9: ELASTICITY
%  The Physics Codex: A Complete University Foundation
%  Korede Samuel Omotosho (Falcon)
%  Aurikrex Learning Systems, 2026
%
%  File: chapters/part1/ch09_elasticity.tex
% ============================================================

\chapter{Elasticity}
\label{ch:elasticity}

\chapterquote{The most incomprehensible thing about
the world is that it is comprehensible.}{Albert Einstein}

\begin{objectivebox}
By the end of this chapter, you should be able to:
\begin{enumerate}[noitemsep]
  \item State and apply Hooke's Law to springs and
  elastic materials
  \item Define elastic limit, yield point and
  breaking point
  \item Define and calculate tensile stress, tensile
  strain and Young's modulus
  \item Interpret force-extension and stress-strain
  graphs
  \item Calculate the energy stored in a stretched
  spring or elastic material
  \item Distinguish between elastic and plastic
  deformation
  \item Classify materials as ductile, brittle or
  elastic and describe their stress-strain curves
  \item Apply spring combinations (series and parallel)
\end{enumerate}
\end{objectivebox}

\vspace{0.4cm}

\minitoc

\vspace{0.5cm}

% ============================================================
\section{Deformation of Solids}
% ============================================================

When a force is applied to a solid object, it deforms
--- it changes shape. Sometimes this deformation is
permanent (like bending a spoon), and sometimes the
object springs back to its original shape when the
force is removed (like a rubber band). Understanding
the relationship between applied forces and the
resulting deformations is essential in engineering,
architecture, and materials science.

\begin{definitionbox}[Elastic and Plastic Deformation]
\textbf{Elastic deformation:} The material returns
to its original shape and size when the deforming
force is removed. The deformation is reversible.

\vspace{0.4cm}

\textbf{Plastic deformation:} The material does
\textit{not} return to its original shape when the
force is removed. The deformation is permanent.
The material has been permanently changed.
\end{definitionbox}

\begin{realworldbox}[Elasticity in Nigerian Construction]
Steel reinforcing bars (``rebar'') used in building
construction across Nigeria must undergo elastic
deformation under normal loads --- if they bent
permanently (plastic deformation) under the weight
of a building, the structure would fail. Engineers
specify steel grades by their yield stress, ensuring
buildings operate safely within the elastic region.
The quality of construction materials is a direct
application of the physics in this chapter.
\end{realworldbox}

% ============================================================
\section{Hooke's Law}
% ============================================================

\begin{lawbox}[Hooke's Law]
Within the \textbf{elastic limit}, the extension
of a spring (or elastic material) is directly
proportional to the applied force:
\[
  F = ke
\]
where:
\begin{itemize}[noitemsep]
  \item $F$ = applied force (N)
  \item $e$ = extension from natural length (m)
  \item $k$ = spring constant (stiffness) (N\,m$^{-1}$)
\end{itemize}
A larger $k$ means a stiffer spring.
\end{lawbox}

\begin{historybox}[Robert Hooke]
Robert Hooke (1635--1703) was one of Newton's great
contemporaries and rivals. He discovered his law of
elasticity in 1660 but kept it secret for years,
publishing it initially as a Latin anagram
(\textit{ceiiinosssttu}) before revealing it as
\textit{Ut tensio, sic vis} --- ``As the extension,
so the force.'' Beyond elasticity, Hooke invented
the universal joint, the iris diaphragm, coined the
term ``cell'' in biology (using his microscope),
and designed many of the scientific instruments of
his era. He is one of the most underrated scientists
in history, largely because his feud with Newton
meant Newton's supporters marginalised him.
\end{historybox}

\begin{diagbox}[Force-Extension Graph (Hooke's Law Region)]
\begin{center}
\begin{tikzpicture}
\begin{axis}[
  width=10cm, height=7cm,
  xlabel={Extension $e$ (m)},
  ylabel={Force $F$ (N)},
  xmin=0, xmax=0.12,
  ymin=0, ymax=70,
  grid=major, grid style={gray!20},
  xtick={0,0.02,0.04,0.06,0.08,0.10,0.12},
  ytick={0,10,20,30,40,50,60,70},
]
  % Hooke's law region (straight line)
  \addplot[codexblue, very thick, domain=0:0.08]{500*x};

  % Beyond elastic limit (curved)
  \addplot[codexred, very thick, domain=0.08:0.115,
    samples=30]{40 + 200*(x-0.08) - 3000*(x-0.08)^2};

  % Elastic limit point
  \fill[codexgold] (axis cs:0.08,40) circle (4pt);
  \draw[dashed, codexgold]
    (axis cs:0.08,0) -- (axis cs:0.08,40);
  \node[codexgold, right, font=\small]
    at (axis cs:0.082,35) {Elastic limit};

  % Gradient label
  \draw[<->, gray]
    (axis cs:0.01,5) -- (axis cs:0.05,25);
  \node[codexblue, font=\small, rotate=63]
    at (axis cs:0.025,20) {gradient $= k$};

  % Labels
  \node[codexblue, font=\small]
    at (axis cs:0.03,38) {\shortstack{Hooke's law\\(linear)}};
  \node[codexred, font=\small]
    at (axis cs:0.10,50) {\shortstack{Beyond elastic\\limit}};
\end{axis}
\end{tikzpicture}
\end{center}
\end{diagbox}

\begin{examplebox}[9.1]
\stepcounter{workedexample}
A spring has natural length $12.0\ \text{cm}$.
When a $4.0\ \text{N}$ weight is hung from it, its
length becomes $17.0\ \text{cm}$.\\
(a) Find the spring constant.\\
(b) Find the total length when a $10.0\ \text{N}$
weight is hung (assuming Hooke's law still holds).\\
(c) Find the elastic PE stored in case (b).
\end{examplebox}

\begin{solutionbox}
\textbf{(a)} Extension $e = 17.0 - 12.0 = 5.0\ \text{cm}
= 0.050\ \text{m}$:
\[
  k = \frac{F}{e} = \frac{4.0}{0.050}
  = 80\ \text{N\,m}^{-1}
\]

\textbf{(b)} Extension for $10\ \text{N}$:
\[
  e = \frac{F}{k} = \frac{10.0}{80} = 0.125\ \text{m}
  = 12.5\ \text{cm}
\]
Total length $= 12.0 + 12.5 = 24.5\ \text{cm}$

\textbf{(c)} Elastic PE:
\[
  E = \frac{1}{2}ke^2
  = \frac{1}{2}(80)(0.125)^2
  = 40 \times 0.015625
  = 0.625\ \text{J}
\]
\end{solutionbox}

% ============================================================
\section{Elastic Potential Energy}
% ============================================================

\begin{processbox}[Deriving Elastic PE]
When a spring is extended by $e$, the force varies
from $0$ (at $e=0$) to $ke$ (at extension $e$).
The work done equals the area under the $F$-$e$ graph:
\[
  W = \text{area of triangle}
  = \frac{1}{2} \times e \times F
  = \frac{1}{2} \times e \times ke
  = \frac{1}{2}ke^2
\]
This work is stored as elastic PE:
\[
  \boxed{E_{elastic} = \frac{1}{2}ke^2
  = \frac{1}{2}Fe = \frac{F^2}{2k}}
\]
\end{processbox}

\begin{examplebox}[9.2]
\stepcounter{workedexample}
A spring of stiffness $200\ \text{N\,m}^{-1}$ is
compressed by $8.0\ \text{cm}$ and used to launch
a $50\ \text{g}$ ball vertically. Assuming all
elastic PE converts to KE, find the launch speed
of the ball.
\end{examplebox}

\begin{solutionbox}
Elastic PE stored:
\[
  E = \frac{1}{2}ke^2
  = \frac{1}{2}(200)(0.08)^2
  = 100 \times 0.0064 = 0.64\ \text{J}
\]

Converting to KE:
\[
  \frac{1}{2}mv^2 = 0.64
\]
\[
  v^2 = \frac{2 \times 0.64}{0.050} = 25.6
\]
\[
  v = \sqrt{25.6} = 5.06\ \text{m\,s}^{-1}
\]
\end{solutionbox}

% ============================================================
\section{Springs in Series and Parallel}
% ============================================================

\begin{lawbox}[Springs in Series and Parallel]
\textbf{Springs in series} (end-to-end):
Both springs experience the same force but the
total extension is the sum of individual extensions:
\[
  \frac{1}{k_{series}} = \frac{1}{k_1} + \frac{1}{k_2}
\]
Series combination is \textit{softer} (smaller $k$,
larger extension for same force).

\vspace{0.5cm}

\textbf{Springs in parallel} (side-by-side):
Both springs have the same extension but the total
force is shared between them:
\[
  k_{parallel} = k_1 + k_2
\]
Parallel combination is \textit{stiffer} (larger $k$,
smaller extension for same force).
\end{lawbox}

\begin{diagbox}[Series vs Parallel Springs]
\begin{center}
\begin{tikzpicture}[scale=0.85]
  % SERIES
  \node[font=\bfseries] at (1,6) {Series};
  % Ceiling
  \fill[gray!30] (0,5.5) rectangle (2,5.8);
  \draw[thick] (0,5.5) -- (2,5.5);
  % Spring 1
  \draw[thick, decorate,
    decoration={coil, amplitude=4pt, segment length=6pt}]
    (1,5.5) -- (1,3.8);
  \node[right, font=\small] at (1.3,4.65) {$k_1$};
  % Spring 2
  \draw[thick, decorate,
    decoration={coil, amplitude=4pt, segment length=6pt}]
    (1,3.8) -- (1,2.1);
  \node[right, font=\small] at (1.3,2.95) {$k_2$};
  % Mass
  \fill[codexblue!40] (0.4,1.7) rectangle (1.6,2.1);
  \node at (1,1.9) {$F$};
  % Formula
  \node[font=\small, align=center] at (1,1.0)
    {$\dfrac{1}{k_s} = \dfrac{1}{k_1}+\dfrac{1}{k_2}$};

  % PARALLEL
  \node[font=\bfseries] at (6,6) {Parallel};
  % Ceiling
  \fill[gray!30] (4,5.5) rectangle (8,5.8);
  \draw[thick] (4,5.5) -- (8,5.5);
  % Spring 1
  \draw[thick, decorate,
    decoration={coil, amplitude=4pt, segment length=6pt}]
    (5,5.5) -- (5,3.0);
  \node[left, font=\small] at (4.7,4.25) {$k_1$};
  % Spring 2
  \draw[thick, decorate,
    decoration={coil, amplitude=4pt, segment length=6pt}]
    (7,5.5) -- (7,3.0);
  \node[right, font=\small] at (7.3,4.25) {$k_2$};
  % Bar connecting springs to mass
  \draw[thick] (4.5,3.0) -- (7.5,3.0);
  % Mass
  \fill[codexblue!40] (5,2.6) rectangle (7,3.0);
  \node at (6,2.8) {$F$};
  % Formula
  \node[font=\small, align=center] at (6,1.0)
    {$k_p = k_1 + k_2$};
\end{tikzpicture}
\end{center}
\end{diagbox}

% ============================================================
\section{Stress, Strain and Young's Modulus}
% ============================================================

Hooke's law $F = ke$ describes the behaviour of a
specific spring. But to characterise the \textit{material}
itself (independent of geometry), we use \textbf{stress}
and \textbf{strain}.

\begin{definitionbox}[Tensile Stress]
\textbf{Tensile stress} $\sigma$ is the force per
unit cross-sectional area applied to a material:
\[
  \sigma = \frac{F}{A}
\]
SI unit: N\,m$^{-2}$ = Pascal (Pa).
A material under tension has positive stress; under
compression, negative.
\end{definitionbox}

\begin{definitionbox}[Tensile Strain]
\textbf{Tensile strain} $\varepsilon$ is the fractional
change in length:
\[
  \varepsilon = \frac{\Delta L}{L_0}
  = \frac{e}{L_0}
\]
Strain is \textbf{dimensionless} (it is a ratio of
two lengths). It is often expressed as a percentage.
\end{definitionbox}

\begin{definitionbox}[Young's Modulus]
The \textbf{Young's modulus} $E$ (or $Y$) of a material
is the ratio of tensile stress to tensile strain within
the elastic limit:
\[
  E = \frac{\sigma}{\varepsilon}
  = \frac{F/A}{\Delta L/L_0}
  = \frac{FL_0}{A\,\Delta L}
\]
SI unit: N\,m$^{-2}$ = Pa (typically GPa for solids).

Young's modulus is a \textit{material property} ---
it does not depend on the size or shape of the
specimen, only on what the material is.
\end{definitionbox}

\begin{table}[H]
\centering
\caption{Young's Modulus of Common Materials}
\label{tab:youngs_modulus}
\begin{tabular}{lll}
\toprule
\textbf{Material} & \textbf{Young's Modulus (GPa)} &
\textbf{Notes}\\
\midrule
Steel & 200 & Construction rebar, bridges\\
Aluminium & 70 & Aircraft, window frames\\
Copper & 120 & Wiring, pipes\\
Concrete & 30 & Compression only; weak in tension\\
Rubber & 0.001--0.1 & Very elastic; large strains\\
Bone & 15--20 & Cortical (compact) bone\\
Wood & 10--15 & Along the grain\\
Glass & 70 & Brittle; fractures suddenly\\
\bottomrule
\end{tabular}
\end{table}

\begin{examplebox}[9.3]
\stepcounter{workedexample}
A steel wire of length $2.0\ \text{m}$ and
cross-sectional area $1.5\times10^{-6}\ \text{m}^2$
is stretched by a $300\ \text{N}$ force.
($E_{steel} = 2.0\times10^{11}\ \text{Pa}$)\\
(a) Calculate the tensile stress.\\
(b) Calculate the extension.\\
(c) Calculate the tensile strain.\\
(d) Calculate the elastic PE stored.
\end{examplebox}

\begin{solutionbox}
\textbf{(a)} Tensile stress:
\[
  \sigma = \frac{F}{A}
  = \frac{300}{1.5\times10^{-6}}
  = 2.0\times10^8\ \text{Pa} = 200\ \text{MPa}
\]

\textbf{(b)} Extension:
\[
  E = \frac{FL_0}{A\Delta L}
  \implies \Delta L = \frac{FL_0}{AE}
  = \frac{300 \times 2.0}{1.5\times10^{-6}
    \times 2.0\times10^{11}}
\]
\[
  E = \frac{600}{3.0\times10^5}
  = 2.0\times10^{-3}\ \text{m} = 2.0\ \text{mm}
\]

\textbf{(c)} Tensile strain:
\[
  \varepsilon = \frac{\Delta L}{L_0}
  = \frac{2.0\times10^{-3}}{2.0}
  = 1.0\times10^{-3} = 0.1\%
\]

\textbf{(d)} Elastic PE:
\[
  E_{pe} = \frac{1}{2}F\Delta L
  = \frac{1}{2} \times 300 \times 2.0\times10^{-3}
  = 0.30\ \text{J}
\]
\end{solutionbox}

% ============================================================
\section{The Stress-Strain Graph}
% ============================================================

The stress-strain graph reveals the complete mechanical
behaviour of a material from initial loading through
to fracture. Each region of the graph has a physical
meaning.

\begin{diagbox}[Stress-Strain Graph for a Ductile Metal (Steel)]
\begin{center}
\begin{tikzpicture}
\begin{axis}[
    width=13cm,
    height=9cm,
    clip=false,
    xlabel={Strain $\varepsilon$},
    ylabel={Stress $\sigma$ (MPa)},
    xmin=0, xmax=0.35,
    ymin=0, ymax=580,
    grid=major,
    grid style={gray!20},
    xtick={0, 0.05, 0.10, 0.15, 0.20, 0.25, 0.30},
    xticklabels={0, 0.05, 0.10, 0.15, 0.20, 0.25, 0.30},
    xticklabel style={rotate=45, anchor=north east, font=\small},
    extra x ticks={0.020},
    extra x tick labels={$\varepsilon_A$},
    extra x tick style={
        tick label style={color=codexblue, font=\small, yshift=-1pt},
    },
    ytick={0,100,200,300,355,400,510},
    yticklabels={$0$,$100$,$200$,$300$,$355$,$400$,$510$},
    tick label style={font=\small},
]

% ===== SHADED ZONES (drawn first, behind everything) =====
\fill[codexblue,  opacity=0.08] (axis cs:0,0)     rectangle (axis cs:0.020,580);
\fill[codexred,   opacity=0.08] (axis cs:0.020,0) rectangle (axis cs:0.060,580);
\fill[codexgreen, opacity=0.08] (axis cs:0.060,0) rectangle (axis cs:0.250,580);
\fill[codexgold,  opacity=0.08] (axis cs:0.250,0) rectangle (axis cs:0.320,580);

% ===== YIELD STRESS REFERENCE LINE =====
\draw[gray, dashed, thin] (axis cs:0,355) -- (axis cs:0.320,355);
\node[gray, font=\scriptsize, anchor=east] at (axis cs:-0.010,355) {$\sigma_y$};

% ===== ZONE LABELS =====
\node[align=center, font=\scriptsize] at (axis cs:0.010,500) {\textbf{Elastic}\\[1pt]reversible};
\node[align=center, font=\scriptsize] at (axis cs:0.040,500) {\textbf{Yield}\\[1pt]plateau};
\node[align=center, font=\scriptsize] at (axis cs:0.155,500) {\textbf{Hardening}\\[1pt]strengthening};
\node[align=center, font=\scriptsize] at (axis cs:0.285,500) {\textbf{Necking}\\[1pt]thinning};

% ===== CURVE SEGMENTS =====
% O -> A : Elastic region (exaggerated scale so it's visible; Hooke's Law)
\addplot[codexblue, very thick, domain=0:0.020, samples=50] {20000*x};

% A -> B' : Drop from upper yield to lower yield
\addplot[codexred, very thick, domain=0.020:0.025, samples=50] {400 - 9000*(x - 0.020)};

% B' -> B : Yield plateau (Luders band)
\addplot[codexred, very thick, domain=0.025:0.060, samples=50] {355};

% B -> C : Strain hardening, eased to zero slope at both ends (no kink at C)
\addplot[codexgreen, very thick, smooth, domain=0.060:0.250, samples=80]
    {355 + 155*(1 - cos(180*(x - 0.060)/0.190))/2};

% C -> D : Necking, eased to zero slope at C (continuous with hardening curve)
\addplot[codexgold, very thick, smooth, domain=0.250:0.320, samples=80]
    {510 - 170*(1 - cos(180*(x - 0.250)/0.070))/2};

% ===== HOOKE'S LAW TRIANGLE (rise/run on the elastic line) =====
\draw[codexblue, thin] (axis cs:0.006,120) -- (axis cs:0.013,120) -- (axis cs:0.013,260);
\node[codexblue, font=\scriptsize, below] at (axis cs:0.0095,118) {$\Delta\varepsilon$};
\node[codexblue, font=\scriptsize, right] at (axis cs:0.0135,190) {$\Delta\sigma$};
\node[codexblue, font=\scriptsize, rotate=80] at (axis cs:0.0165,300) {$E=\Delta\sigma/\Delta\varepsilon$};
\node[codexblue, font=\scriptsize, rotate=80, anchor=east] at (axis cs:0.0005,250) {Hooke's Law region};

% Not-to-scale disclaimer
\node[gray, font=\tiny] at (axis cs:0.075,55) {(elastic region exaggerated for clarity --- not to scale)};

% ===== POINT MARKERS =====
\fill (axis cs:0,0) circle (1.6pt);
\node[below right, font=\small] at (axis cs:0,0) {O\,(origin)};

\fill[codexblue] (axis cs:0.020,400) circle (1.6pt);
\node[above right, codexblue, font=\small] at (axis cs:0.022,405) {A\,(upper yield)};

\fill[codexred] (axis cs:0.060,355) circle (1.6pt);
\node[below right, codexred, font=\small] at (axis cs:0.060,355) {B\,(lower yield)};

\fill[codexgreen] (axis cs:0.250,510) circle (1.6pt);
\node[above, codexgreen, font=\small] at (axis cs:0.250,515) {C\,(UTS)};

\fill[codexgold] (axis cs:0.320,340) circle (1.6pt);
\node[above right, codexgold, font=\small] at (axis cs:0.320,340) {D\,(fracture)};

\end{axis}
\end{tikzpicture}
\end{center}

\vspace{0.4cm}
{\small
\renewcommand{\arraystretch}{1.5}
\begin{center}
\begin{tabular}{p{1.6cm} | p{10.5cm}}
\hline
\textbf{Region} & \textbf{Description} \\
\hline
O $\to$ A & \textbf{Elastic region} --- Hooke's Law obeyed; slope $=$ Young's modulus $E$ \\
A $\to$ B & \textbf{Yield plateau} --- large strain for little stress increase; material flows (L\"{u}ders band in mild steel) \\
B $\to$ C & \textbf{Strain hardening} --- dislocations accumulate; stress required increases \\
C & \textbf{Ultimate Tensile Strength (UTS)} --- maximum engineering stress \\
C $\to$ D & \textbf{Necking} --- cross-section reduces locally; engineering stress falls \\
D & \textbf{Fracture point} \\
\hline
\end{tabular}
\end{center}
}
\end{diagbox}

\begin{defbox}[Key Points on the Stress-Strain Graph]
\begin{itemize}[noitemsep]
  \item \textbf{O to A (Elastic region):} Stress is
  proportional to strain. Hooke's Law is obeyed.
  Young's modulus $= $ gradient.
  \item \textbf{A (Elastic limit / Yield point):}
  Maximum stress for which the material returns to
  original shape. Beyond this, deformation is permanent.
  \item \textbf{A to B (Plastic deformation / Yield
  plateau):} Large strain for little increase in stress.
  Material flows.
  \item \textbf{B to C (Strain hardening):} Material
  becomes harder to deform as dislocations pile up.
  Stress required increases.
  \item \textbf{C (Ultimate Tensile Strength, UTS):}
  Maximum stress the material can withstand.
  \item \textbf{C to D (Necking):} Material becomes
  locally thinner (necking). Stress in that region
  increases dramatically until fracture.
  \item \textbf{D (Fracture point):} Material breaks.
\end{itemize}
\end{defbox}

\subsection{Material Classifications}

\begin{defbox}[Ductile, Brittle and Elastic Materials]
\df{Ductile}{Large plastic region before fracture.
Can be drawn into wires or hammered into sheets.
Examples: copper, steel, gold, aluminium.}

\vspace{0.4cm}

\df{Brittle}{Little or no plastic region.
Fractures suddenly with little warning when the elastic
limit is exceeded. Examples: glass, cast iron, ceramics,
chalk, concrete.}

\vspace{0.4cm}

\df{Elastic (rubber-like)}{Large elastic strains,
non-linear stress-strain relationship (does not obey
Hooke's Law over large ranges). High resilience.
Examples: rubber, biological tissues.}
\end{defbox}

\begin{diagbox}[Comparing Stress-Strain Curves]
\begin{center}
\begin{tikzpicture}
\begin{axis}[
  width=10cm, height=6cm,
  xlabel={Strain $\varepsilon$},
  ylabel={Stress $\sigma$},
  xmin=0, xmax=1.0,
  ymin=0, ymax=1.2,
  grid=major, grid style={gray!20},
  xtick=\empty, ytick=\empty,
  legend pos=north west
]
  % ---- Ductile (steel-like): two segments, ONE legend entry ----
  \addplot[codexblue, very thick, domain=0:0.02,
    samples=10, forget plot]{30*x};
  \addplot[codexblue, very thick, domain=0.02:0.4,
    samples=20]{0.6 + 0.6*(x-0.02) - 1*(x-0.02)^2};
  \addlegendentry{Ductile (steel)}
  \fill[codexblue] (axis cs:0.4,0.684) circle (2pt); % fracture point

  % ---- Brittle (glass-like): straight line then sudden fracture ----
  \addplot[codexred, very thick, domain=0:0.04,
    samples=10]{15*x};
  \addlegendentry{Brittle (glass)}
  \fill[codexred] (axis cs:0.04,0.6) circle (2pt);
  \draw[codexred, very thick, dashed]
    (axis cs:0.04,0.6) -- (axis cs:0.04,0);

  % ---- Elastic/rubber: J-shaped curve, low initial stiffness then steep upturn ----
  \addplot[codexgreen, very thick, domain=0:0.95,
    samples=60]{0.12*x + 1.0*x^4};
  \addlegendentry{Elastic (rubber)}
\end{axis}
\end{tikzpicture}
\end{center}
\end{diagbox}

% ============================================================
\section{Chapter Summary}
% ============================================================

\begin{summarybox}
\textbf{Hooke's Law:} $F = ke$ (within elastic limit).
Spring constant $k$ (N\,m$^{-1}$) measures stiffness.

\textbf{Elastic PE:}
$E = \frac{1}{2}ke^2 = \frac{1}{2}Fe = \frac{F^2}{2k}$

\textbf{Springs in series:}
$\frac{1}{k_s} = \frac{1}{k_1} + \frac{1}{k_2}$
(softer)

\textbf{Springs in parallel:}
$k_p = k_1 + k_2$ (stiffer)

\textbf{Tensile stress:}
$\sigma = F/A$ (Pa)

\textbf{Tensile strain:}
$\varepsilon = \Delta L/L_0$ (dimensionless)

\textbf{Young's modulus:}
$E = \sigma/\varepsilon = FL_0/(A\Delta L)$ (Pa)

\textbf{Elastic deformation:} reversible.
\textbf{Plastic deformation:} permanent.
\textbf{Elastic limit:} boundary between them.

\textbf{Material types:}
Ductile (large plastic region);
Brittle (fractures suddenly at elastic limit);
Elastic/rubber (large, non-linear elastic strains).
\end{summarybox}

% ============================================================
\newpage
\begin{exercisebox}[9]

\subsection*{Section A: Multiple Choice}

\textbf{A1.} A spring of constant $k = 40\ \text{N\,m}^{-1}$
is stretched by $5\ \text{cm}$. The elastic PE stored is:

\mcoption{A}{$0.025\ \text{J}$}
\mcoption{B}{$0.05\ \text{J}$}
\mcoption{C}{$0.10\ \text{J}$}
\mcoption{D}{$1.00\ \text{J}$}
\ansref{Ch9-A1}

\vspace{0.4cm}

\textbf{A2.} Two springs each of constant
$k = 60\ \text{N\,m}^{-1}$ are connected in series.
The effective spring constant is:

\mcoption{A}{$120\ \text{N\,m}^{-1}$}
\mcoption{B}{$60\ \text{N\,m}^{-1}$}
\mcoption{C}{$30\ \text{N\,m}^{-1}$}
\mcoption{D}{$15\ \text{N\,m}^{-1}$}
\ansref{Ch9-A2}

\vspace{0.4cm}

\textbf{A3.} A wire of length $2\ \text{m}$ and
cross-section $2\ \text{mm}^2$ is stretched by
$0.5\ \text{mm}$ under a $100\ \text{N}$ load.
The Young's modulus is:

\mcoption{A}{$2\times10^{10}\ \text{Pa}$}
\mcoption{B}{$2\times10^{11}\ \text{Pa}$}
\mcoption{C}{$1\times10^{11}\ \text{Pa}$}
\mcoption{D}{$5\times10^{10}\ \text{Pa}$}
\ansref{Ch9-A3}

\vspace{0.4cm}

\textbf{A4.} Which material behaviour is described
as brittle?

\mcoption{A}{Large plastic deformation before fracture}
\mcoption{B}{Returns to original shape after large
deformation}
\mcoption{C}{Fractures with little or no plastic
deformation}
\mcoption{D}{Has a very small Young's modulus}
\ansref{Ch9-A4}

\vspace{0.4cm}

\textbf{A5.} On a stress-strain graph, the gradient
of the linear (elastic) region equals:

\mcoption{A}{The spring constant}
\mcoption{B}{The elastic limit}
\mcoption{C}{The Young's modulus}
\mcoption{D}{The ultimate tensile strength}
\ansref{Ch9-A5}

\vspace{0.4cm}

\textbf{A6.} A rubber band is stretched to twice its
natural length. The strain is:

\mcoption{A}{$0.5$}
\mcoption{B}{$1.0$}
\mcoption{C}{$2.0$}
\mcoption{D}{$100\%$ which equals 1.0}
\ansref{Ch9-A6}

\vspace{0.4cm}

\textbf{A7.} If the diameter of a wire is doubled
(keeping length and material the same), and the same
force is applied, the extension:

\mcoption{A}{Doubles}
\mcoption{B}{Halves}
\mcoption{C}{Becomes one quarter}
\mcoption{D}{Is unchanged}
\ansref{Ch9-A7}

\vspace{0.4cm}

\textbf{A8.} Beyond the elastic limit of a metal wire:

\mcoption{A}{The wire returns to its original
length when the force is removed}
\mcoption{B}{The wire breaks immediately}
\mcoption{C}{Permanent deformation occurs}
\mcoption{D}{Hooke's Law is still obeyed}
\ansref{Ch9-A8}

\vspace{0.4cm}

\textbf{A9.} A spring of constant $k$ stores energy
$E$ when extended by $x$. If the extension is doubled,
the stored energy becomes:

\mcoption{A}{$2E$}
\mcoption{B}{$4E$}
\mcoption{C}{$E\sqrt{2}$}
\mcoption{D}{$E/2$}
\ansref{Ch9-A9}

\vspace{0.4cm}

\textbf{A10.} Three springs each of constant
$k = 90\ \text{N\,m}^{-1}$ are connected in parallel.
The effective spring constant is:

\mcoption{A}{$30\ \text{N\,m}^{-1}$}
\mcoption{B}{$90\ \text{N\,m}^{-1}$}
\mcoption{C}{$180\ \text{N\,m}^{-1}$}
\mcoption{D}{$270\ \text{N\,m}^{-1}$}
\ansref{Ch9-A10}

\vspace{0.6cm}

\subsection*{Section B: Theory and Calculations}

\textbf{B1.} A student sets up an experiment to
verify Hooke's Law using a spring. She records the
following data:

\begin{center}
\begin{tabular}{cc}
\toprule
Force (N) & Length (cm)\\
\midrule
0.0 & 15.0\\
2.0 & 17.4\\
4.0 & 19.8\\
6.0 & 22.2\\
8.0 & 24.6\\
10.0 & 26.9\\
12.0 & 30.1\\
\bottomrule
\end{tabular}
\end{center}

(a) Plot force against extension and determine the
spring constant from the gradient.\\
(b) Identify the approximate elastic limit from
the data.\\
(c) What physical change occurs in the spring at
the elastic limit?\\
(d) Calculate the elastic PE stored when the force
is $8.0\ \text{N}$.
\hfill\ansref{Ch9-B1}

\vspace{0.4cm}

\textbf{B2.} A steel cable in a suspension bridge
has the following properties:
Length $= 500\ \text{m}$, diameter $= 50\ \text{mm}$,
Young's modulus $= 200\ \text{GPa}$.
The cable supports a load of $2.0\times10^6\ \text{N}$.\\[4pt]
(a) Calculate the cross-sectional area.\\
(b) Calculate the tensile stress.\\
(c) Calculate the extension of the cable.\\
(d) Calculate the tensile strain and express as a
percentage.\\
(e) Calculate the elastic PE stored in the cable.
\hfill\ansref{Ch9-B2}

\vspace{0.4cm}

\textbf{B3.} A spring of constant
$500\ \text{N\,m}^{-1}$ and natural length
$20\ \text{cm}$ hangs vertically. A $2.0\ \text{kg}$
mass is attached and released from rest at the
natural length position.\\[4pt]
(a) Find the equilibrium extension (where the
net force is zero).\\
(b) Find the amplitude of oscillation.
(Hint: the system oscillates about the equilibrium
position, and the amplitude equals the equilibrium
extension for a release from the natural length.)\\
(c) Find the maximum speed of the mass.\\
(d) Find the maximum elastic PE and maximum KE
during the oscillation.\\
(e) Verify energy conservation by showing that total
mechanical energy is constant at equilibrium and
at maximum extension.
\hfill\ansref{Ch9-B3}

\vspace{0.4cm}

\textbf{B4.} Two springs, $k_1 = 40\ \text{N\,m}^{-1}$
and $k_2 = 60\ \text{N\,m}^{-1}$, are connected in
series and hung from a fixed support. A $3.0\ \text{kg}$
mass hangs from the lower spring.
($g = 10\ \text{m\,s}^{-2}$)\\[4pt]
(a) Calculate the effective spring constant.\\
(b) Calculate the total extension.\\
(c) Calculate the individual extensions of each
spring.\\
(d) Verify that the force in each spring is the same.\\
(e) The mass is now connected to two springs in
parallel (both $k_1$ and $k_2$ support the same mass).
Find the new effective spring constant and total
extension.
\hfill\ansref{Ch9-B4}

\vspace{0.4cm}

\textbf{B5.} A copper wire and a steel wire are each
$1.0\ \text{m}$ long and $0.5\ \text{mm}$ in diameter.
They are joined end-to-end (in series) and a force
of $50\ \text{N}$ is applied.
($E_{copper} = 1.2\times10^{11}\ \text{Pa}$,
$E_{steel} = 2.0\times10^{11}\ \text{Pa}$)\\[4pt]
(a) Calculate the cross-sectional area of each wire.\\
(b) Calculate the extension of each wire.\\
(c) Calculate the total extension of the combined wire.\\
(d) Calculate the stress in each wire. Are they the
same? Explain why.\\
(e) Calculate the strain in each wire. Are they the
same? Explain why.
\hfill\ansref{Ch9-B5}

\end{exercisebox}